%%%%%%%%%%%%%%%%%%%%%%%%%%%%%%%%%%%%%%%%%%%%%%%%%%%%
% Template para utilização do Beamer
% Apresentação para Universidade Federal de Catalão
% Versão 1.0 (03/03/2023)
%%%%%%%%%%%%%%%%%%%%%%%%%%%%%%%%%%%%%%%%%%%%%%%%%%%%

\documentclass{beamer}

\mode<presentation> {

%------------------------------------------------
% A classe Beamer possui vários temas de slides.
% A lista a seguir contém os temas.
% Remova o símbolo da porcentagem para utilização.

%\usetheme{default}
%\usetheme{AnnArbor}
%\usetheme{Antibes}
%\usetheme{Bergen}
%\usetheme{Berkeley}
%\usetheme{Berlin}
%\usetheme{Boadilla}
%\usetheme{CambridgeUS}
%\usetheme{Copenhagen}
%\usetheme{Darmstadt}
%\usetheme{Dresden}
%\usetheme{Frankfurt}
%\usetheme{Goettingen}
%\usetheme{Hannover}
%\usetheme{Ilmenau}
%\usetheme{JuanLesPins}
%\usetheme{Luebeck}
\usetheme{Madrid}
%\usetheme{Malmoe}
%\usetheme{Marburg}
%\usetheme{Montpellier}
%\usetheme{PaloAlto}
%\usetheme{Pittsburgh}
%\usetheme{Rochester}
%\usetheme{Singapore}
%\usetheme{Szeged}
%\usetheme{Warsaw}
%------------------------------------------------

%------------------------------------------------
% A classe Beamer possui vários temas de cores.
% Pode ser utilizado em qualquer tema de slide escolhido.

%\usecolortheme{albatross}
%\usecolortheme{beaver}
%\usecolortheme{beetle}
%\usecolortheme{crane}
%\usecolortheme{dolphin}
%\usecolortheme{dove}
%\usecolortheme{fly}
%\usecolortheme{lily}
%\usecolortheme{orchid}
%\usecolortheme{rose}
%\usecolortheme{seagull}
%\usecolortheme{seahorse}
%\usecolortheme{whale}
%\usecolortheme{wolverine}
%------------------------------------------------

%------------------------------------------------
% Remover a linha do rodapé em todos os slides.

%\setbeamertemplate{footline}
%------------------------------------------------

%------------------------------------------------
% Substituir a linha de rodapé em todos os slides por uma contagem simples.

%\setbeamertemplate{footline}[page number]
%------------------------------------------------

%------------------------------------------------
% Remover os símbolos de navegação da parte inferior de todos os slides.

%\setbeamertemplate{navigation symbols}{}
%------------------------------------------------
}

\usepackage{graphicx} % Incluir imagens.
\usepackage{booktabs} % Uso de \toprule, \midrule e \bottomrule em tabelas.

%------------------------------------------------
% Primeiro Slide
\title[Subtitulo aqui]{Titulo aqui}

\author{Nome aqui}
\institute[Sigla opcional]
{
Nome da instituição aqui \\
\medskip
\textit{email aqui}
}
\date{}

\logo{\includegraphics[height=0.8cm]{Logo_UFCAT.png}}

\begin{document}

\begin{frame}
\titlepage
\end{frame}
%------------------------------------------------

%------------------------------------------------

\begin{frame}
\frametitle{Slide 1}

\begin{block}{}
Bloco de texto, equação ou lista 1
\end{block}

\begin{block}{}
Bloco de texto, equação ou lista 2
\end{block}

\end{frame}
%------------------------------------------------

\begin{frame}
\frametitle{Slide 2}

\begin{itemize}

\item Lista de itens 1

\begin{itemize}
    \item Item 1
    \item Item 2
    \item Item 3
\end{itemize}

\item Lista de itens 2

\begin{itemize}
    \item Item 1
\end{itemize}

\item Item 1
\item Item 2
\item Item 3

\end{itemize}

\end{frame}
%------------------------------------------------

\begin{frame}
\frametitle{Slide 3}

\begin{equation}
	Equacao 1
	\label{citacao1}
\end{equation}

\begin{equation}
    Equacao 2
    \label{citacao2}
\end{equation}

\end{frame}
%------------------------------------------------

\begin{frame}
\frametitle{Slide 4}

\begin{block}{Bloco de texto com título 1}
Bloco de texto, equação ou lista 1.
\end{block}

\begin{block}{Bloco de texto com título 2}
Bloco de texto, equação ou lista 2.
\end{block}

\end{frame}
%------------------------------------------------

\begin{frame}
\frametitle{Slide 5}

Inserção de imagem.
\begin{figure}[!h]
	\centering
    %\caption{Título aqui}
	%\includegraphics[width=0.8\linewidth]{Caminho da imagem aqui}
    %\label{Citação aqui}
\end{figure}

\end{frame}
%------------------------------------------------

\begin{frame}
\frametitle{Slide 6}

\begin{block}{Bloco com título contendo exemplo de sistema}
$
y=\left\{\begin{array}{ccc}
	+ \gamma & se & x > 0\\
	- \gamma & se & x \leq 0.
\end{array}\right.
$
\end{block}
\end{frame}
%------------------------------------------------

\begin{frame}
\frametitle{Slide 7}

\begin{table}[!h]
\center
\caption{Exemplo de tabela.}
\begin{tabular}{|c|c|}
\hline 
Item 1 & Descrição ou valor do item 1. \\ 
\hline 
Item 2 & Descrição ou valor do item 2. \\ 
\hline 
\end{tabular} 
\end{table}

\end{frame}
%------------------------------------------------

\begin{frame}
\frametitle{Slide 8}

Inserção de uma imagem do lado da outra.
\begin{figure}[!h]
    \centering
    \begin{minipage}{0.5\textwidth}
        \centering
        %\includegraphics[width=0.7\linewidth]{Caminho da imagem 1 aqui} 
    \end{minipage}\hfill
    \begin{minipage}{0.5\textwidth}
        \centering
        %\includegraphics[width=0.6\linewidth]{Caminho da imagem 2 aqui}
    \end{minipage}
\end{figure}

\end{frame}
%------------------------------------------------

\begin{frame}
\frametitle{Slide 9}

Inserção de várias imagens.
\begin{figure}[!h]
    \centering
    %\includegraphics[width=0.3\linewidth]{Caminho da imagem 1 aqui}
    %\includegraphics[width=0.3\linewidth]{Caminho da imagem 2 aqui}
    %\includegraphics[width=0.3\linewidth]{Caminho da imagem 3 aqui}
\end{figure}

\end{frame}
%----------------------------------------------

\end{document}
